\section{先行関係と並行作業：機械と部品の二つの順序}
\label{theory:discrete:precedence}
複数の機械で複数の部品を加工するとき、同じ機械を同時に使わない資源条件と、同じ部品の加工順の条件を別々に書く。機械AとBが部品p,q,rを各p→q→r順に処理し、各部品がA→Bへ進むなら、A列の順、B列の順、各部品のAからBへの順が全て必要である。

アロー図の矢印は作業、節点はその完了・開始の条件を表す。ある節点へ複数の矢印が入るなら、その全てが完了してから次へ進む。ダミー作業は実際の加工を増やすのでなく、必要な先行関係を追加する。例えばAqとBpの両方が終わってからBqを始めたい場合、Bpの終了節点へAqの終了からダミーを入れればよい。

候補の割当ては、見た目の左右で判断せず、各機械の順序と各部品の順序を一つずつ照合する。不要な待ちを追加するだけの場合と、必要な順を保証しない場合も異なる。並行できるという記述を全作業同時開始と考えず、共有する機械・未完了の前工程の有無を調べる。
